\begin{frame}{1D 단면: 걸음 크기가 속도를 정한다}
  \begin{columns}[T]
    \begin{column}{0.46\textwidth}
      \centering
      \includegraphics[width=\linewidth]{ch02_gradient_descent_1d.png}
      \vskip0.4em
      {\footnotesize 비용 $J(w_1)$ 의 1차원 단면 포물선 위에 세
      학습률($\alpha = 0.1, 0.01, 0.001$)의 경로 --- 같은 출발점.}
    \end{column}
    \begin{column}{0.52\textwidth}
      \begin{itemize}
        \item 세 곡선은 \textbf{같은 포물선 위의 같은 점}에서 출발.
        \item $\alpha = 0.1$ 이 가장 빠르게 바닥에 닿고,
          $\alpha = 0.01$ 은 약 \textbf{10배 느리다}.
        \item $\alpha = 0.001$ 은 300스텝이 지나도 바닥에서 꽤 먼
          채로 거의 움직이지 못함.
        \item 1D 그림은 ``속도''를 보여준다 --- 학습률은
          ``가는 방향''이 아니라 \textbf{``걸음 수''}를 정하는
          하이퍼파라미터.
      \end{itemize}
    \end{column}
  \end{columns}
\end{frame}
